\begin{lemma}[Finite scalar normalization for thinning]
Suppose the extended nonnegative scalars satisfy
$0<\beta<\eta$, $0<\rho<D\leq1$, $\xi>0$,
$4((\rho/D)\beta+\xi)\leq\eta$, and $10\eta<\rho$.
Then $D,\beta,\rho,\eta,\xi$ are all finite, $\eta<\rho$, $\beta\leq D$,
and the real conversions of $\eta,\xi$ are positive. There is a
nonnegative real $q\leq1$ such that
$qD.\operatorname{toReal}=\rho.\operatorname{toReal}$ and
$q\beta.\operatorname{toReal}+\xi.\operatorname{toReal}
\leq\eta.\operatorname{toReal}$.
\end{lemma}
\begin{proof}
Since $\eta\leq10\eta<\rho<D\leq1$, these three scalars and $D$ are
finite, and $\eta<\rho$. The inequality $\beta<\eta$ gives finiteness
of $\beta$ and $\beta<D$. Nonnegativity gives
$\xi\leq(\rho/D)\beta+\xi\leq4((\rho/D)\beta+\xi)\leq\eta$, so $\xi$
is also finite. Positive finite extended nonnegative reals have positive
real conversions, establishing the asserted positivity and also
$D.\operatorname{toReal}>0$.
Set $q=\rho.\operatorname{toReal}/D.\operatorname{toReal}$. Finite real
conversion preserves $\rho<D$, hence $0\leq q\leq1$, and multiplication
by the positive denominator proves $qD.\operatorname{toReal}
=\rho.\operatorname{toReal}$. The quotient $\rho/D$ is finite because
its numerator is finite and its denominator is positive. Thus every
sum and product in the budget is finite, and converting that budget
to real numbers gives $4(q\beta.\operatorname{toReal}
+\xi.\operatorname{toReal})\leq\eta.\operatorname{toReal}$.
The parenthesized quantity is nonnegative; removing the factor four
proves the required error inequality.
\end{proof}
